\documentclass[10pt,a4paper]{amsart}
\usepackage[margin=19mm]{geometry}
\usepackage{amsmath,amssymb,mathtools}
\usepackage[scr=rsfs,cal=euler]{mathalfa}
\usepackage{xfrac}
\usepackage[unicode,hidelinks,bookmarks=false]{hyperref}
\newcommand{\ie}{i.e.\ }
\newcommand{\cf}{cf.\ }
\newcommand{\resp}{respectively\ }
\newcommand{\Subsection}{\subsection}
\providecommand{\nobreakdash}{\nobreak\hbox{-}}
\setlength{\emergencystretch}{2em}
\multlinegap0pt
\usepackage{bm}
\usepackage{amssymb}
\usepackage{xcolor}
\usepackage{algorithm}\usepackage{algpseudocode}
\newtheorem{corollary}{Corollary}
\newtheorem{lemma}{Lemma}
\newcommand{\E}{\mathbb{E}}
\newcommand{\R}{\mathbb{R}}
\title{Uniform Scaling Limits in AdamW-Trained Transformers}
\author{William Gibson, Christoph Reisinger}
\date{Source: arXiv:2605.11059v1 (CC BY 4.0)}
\begin{document}
\maketitle

\noindent\emph{Source and licence.} Uniform Scaling Limits in AdamW-Trained Transformers. Version arXiv:2605.11059v1 (CC BY 4.0), distributed by arXiv under the Creative Commons Attribution 4.0 International licence, http://creativecommons.org/licenses/by/4.0/, as recorded in the arXiv licence metadata for this exact version and shown on the abstract page https://arxiv.org/abs/2605.11059v1. Full licence notice: \texttt{licenses/arxiv-2605.11059-CC-BY-4.0.txt}.\par\smallskip

\noindent\textit{Editorial scope.} Counted unit: the article's corollary cor:LipGammaDerivs (source0.tex lines 647--653). The article's complete original proof of it is delivered with it (source0.tex lines 654--712, one \texttt{proof} environment of 59 lines). No mathematics is written, repaired or re-derived: the statement and the proof are the article's own lines, in the article's own notation, and no retained line is substituted or re-flowed.  Retained uncounted as the source-local context the proof needs: lines 94--97; lines 158--165; lines 414--420; lines 611--617; lines 634--636.  Nothing was omitted from any retained span, so the package carries no omission record.  No retained line contains a citation, so the wrapper carries no bibliography and no bibliography entry.  No retained line contains a figure environment, an image-inclusion command or drawing code, so no image is delivered and the package declares no asset dependency.  Editorial formatting only, in addition: an AMS \texttt{amsart} preamble that restates the article's own declarations and macros used by the retained lines, namely the environments \texttt{corollary}, \texttt{lemma} on separate counters, together with the standard packages those lines need.  The retained environments print in this document as Corollary 1, Lemma 1 in order of appearance, whereas the article prints the same items with the numbering of its own full document.  The complete native source of the article as distributed in the arXiv e-print for this exact version is bundled unchanged as \texttt{sources/200337/source0.tex}.
\begin{equation}\label{eqn:gamma}
\gamma : \R^d\times \mathcal P_c(\R^d)\to \R^d, \qquad
    \gamma(z,\mu) = \frac{\int \exp( \left\langle z,  y \right\rangle)\,y\,d\mu(y)}{\int \exp( \left\langle z, y \right\rangle)\,d\mu(y)},
\end{equation}

\begin{corollary}\label{corr:LipMu} 
Take any $z \in \mathbb{R}^d$ and $p\in [1,\infty]$. Then, for any $\mu,\tilde{\mu}\in \mathcal{P}_p(\bar B(R))$, there exists $\Lambda_{\mu,p}: \mathbb R_{\ge0}\to \mathbb R_{\ge 0}$  such that
    \[ \Vert \gamma^R(z,\mu)-\gamma^R(z,\tilde{\mu})\Vert \le \Lambda_{\mu,p}(\Vert z \Vert) W_p(\mu, \tilde{\mu}), \]
    where $\Lambda_{\mu,p}$ is given by
$
\Lambda_{\mu,p}(\tilde R) =(1+2R\tilde R)\exp(2R\tilde Rp^{-1})
$.
\end{corollary}

\begin{lemma}\label{ProbMeasureMVT}
    Suppose $f:\mathcal{P}_2(\mathbb{R}^d) \mapsto \mathbb{R}$ is a continuous function such that its Lions derivative $\partial_{\mu}f: \mathcal{P}_2(\mathbb{R}^d)\times \mathbb{R}^d \mapsto \mathbb{R}^d$ exists. Then there exists $t\in (0,1)$ such that 
    \[
    f(\tilde\mu)-f(\mu) = \widetilde{\mathbb{E}}\left[\left\langle\partial_{\mu}f\left(\mathcal{L}(Z + t(\tilde{Z}-Z)),Z + t(\tilde{Z}-Z)\right),\tilde{Z}-Z\right\rangle\right],
    \]
    for any $\mu,\tilde{\mu}\in \mathcal{P}_2(\mathbb{R}^d)$ and random variables $Z,\tilde{Z}$, defined on an atomless Polish probability space $(\widetilde{\Omega}, \widetilde{\mathcal{F}},\widetilde{\mathbb{P}})$, such that $\mathcal{L}(Z)=\mu,\mathcal{L}(\tilde{Z})=\tilde{\mu}$. Here, $\widetilde{\E}$ is the expectation with respect to the measure $\widetilde{\mathbb P}$.
\end{lemma}

\begin{lemma}\label{lemma:2ndderivatives}
Given $\gamma: \R^d\times \mathcal P_2(\R^d)\to \R^d$ defined in \eqref{eqn:gamma} and radius $R_1>0$, for any\\ $y\in \bar B(R_1), \mu \in\mathcal P(\bar B(R_1))$ and $z,\Delta z\in \R^d$, we have
    \begin{align*}
                \left\Vert D_{zz}\gamma(z,\mu)\right\Vert_{\mathrm{HS}} &\le 8 R_1^3, \\ 
        \left\Vert \sum_{k=1}^d\partial_{z_k}\partial_\mu\gamma(z,\mu)(y)(\Delta z)_k\right\Vert_\mathrm{op}&\le 2R_1(2+3R_1\Vert z\Vert)\frac{\exp(z^Ty)}{\gamma_2(z,\mu)}\Vert \Delta z \Vert.
    \end{align*}
\end{lemma}

 \begin{equation}\label{eqn:Dmugamma}
               \partial_\mu \gamma(z,\mu)(y) = \frac{\exp(z^Ty)}{\gamma_2(z,\mu)}\left[  \{y-\gamma(z,\mu)\} z^T+I\right]. 
 \end{equation}

\begin{corollary}\label{cor:LipGammaDerivs}
Given $R_1,R_z>0$, there exists a constant $\Lambda_{D\gamma,p}=\Lambda_{D\gamma,p}(R_1,R_z)$ such that for any $y_i\in \bar B(R_1), z_i\in \bar B(R_z)$ and $\mathfrak{m}_i\in \mathcal P(\bar B(R_1))$ for $i=1,2$, we have that
\begin{align*}
    \Vert D_z\gamma(z_1,\mathfrak{m}_1)-D_z\gamma(z_2,\mathfrak{m}_2)\Vert_\mathrm{op}&\le \Lambda_{D\gamma,p} \left(\Vert z_1-z_2\Vert +W_p(\mathfrak{m}_1,\mathfrak{m}_2) \right),\\
    \Vert \partial_\mu\gamma(z_1,\mathfrak{m}_1)(y_2)-\partial_\mu\gamma(z_2,\mathfrak{m}_2)(y_2)\Vert_\mathrm{op}&\le\Lambda_{D\gamma,p} \left(\Vert y_1-y_2\Vert+\Vert z_1-z_2\Vert +W_p(\mathfrak{m}_1,\mathfrak{m}_2) \right).
\end{align*}
\end{corollary}

\begin{proof}
Let $(\widetilde{\Omega},\widetilde{\mathcal F},\widetilde{\mathbb P})$ be an atomless Polish probability space. Then, for any $u,v \in \mathbb S^{d-1}$, $t\in(0,1)$ and $Y_i\in L^2(\widetilde{\Omega},\widetilde{\mathcal F},\widetilde{\mathbb P})$ with $\mathcal L(Y_i)=\mathfrak{m}_i$ for $i=1,2$, the Mean Value Theorem given in Lemma \ref{ProbMeasureMVT} implies that 
\[
u^T\left(D_{z}\gamma(z_2,\mathfrak{m}_1)-D_{z}\gamma(z_2,\mathfrak{m}_2)\right)v = \sum_{i=1}^d\widetilde{\mathbb E} \left[u_i\, v^T\partial_{z_i}\partial_\mu \gamma(z_2,\mathcal L(Y_t))(Y_t)(Y_1-Y_2) \right] ,
\]
where $Y_t=Y_2+t(Y_1-Y_2)$. Therefore, H\"older's inequality implies that for any $p,q\in [1,\infty]$ such that $p^{-1}+q^{-1}=1$, we have 
\begin{align*}
   u^T\left( D_{z}\gamma(z_2,\mathfrak{m}_1)-D_{z}\gamma(z_2,\mathfrak{m}_2)\right)v = \left\Vert \left\Vert \sum_{i=1}^du_i \partial_{z_i}\partial_\mu \gamma(z_2,\mathcal L(Y_t))(Y_t)\right\Vert_\mathrm{op}\right\Vert_{L^q(\widetilde{\Omega})} \Vert Y_1-Y_2\Vert_{L^p(\widetilde{\Omega})}.
\end{align*}
Hence, by Lemma \ref{lemma:2ndderivatives} and since $\Vert u \Vert =1$,
\begin{align*}
    u^T\left( D_{z}\gamma(z_2,\mathfrak{m}_1)-D_{z}\gamma(z_2,\mathfrak{m}_2)\right)v = 2R_1(2+3R_1R_z)\left\Vert \frac{\exp(z^TY_t)}{\gamma_2\left(z,\mathcal L(Y_t)\right)}\right\Vert_{L^q(\widetilde{\Omega})} \Vert  Y_2-Y_1\Vert_{L^p(\widetilde{\Omega})}.
\end{align*}
Evaluating the above expression at the coupling between $Y_1,Y_2$ that attains the $p$-Wasserstein distance,  we obtain
\begin{equation*}
    u^T\left(D_z\gamma(z_2,\mathfrak{m}_1)-D_{z}\gamma(z_2,\mathfrak{m}_2)\right)v\le 3R_1(1+2R_1R_z) \exp(2R_zR_1p^{-1}) W_p(\mathfrak{m}_1,\mathfrak{m}_2),
\end{equation*}
 where we have used $\Vert Y_t\Vert \le R_1 $ $\widetilde{\mathbb P}$-a.s. and  $1-q^{-1}=p^{-1}$ to produce the bound $$\left\Vert \frac{\exp(z^TY_t)}{\gamma_2(z,\mathcal L(Y_t))}\right\Vert_{L^q(\widetilde{\Omega})}=\left(\widetilde{\mathbb E}\left[\frac{\exp(z^TY_t)}{\widetilde{\mathbb E}\left[\exp(z^TY_t)\right]}\cdot \frac{\exp((q-1)z^TY_t)}{(\gamma_2(z,\mathcal L(Y_t)))^{q-1}}\right]\right)^{1/q}\le \exp\left(2R_1R_zp^{-1}\right).$$ 
 The preceding bound on the difference between $D_z\gamma(z_2,\mathfrak{m}_1)$ and $D_z\gamma(z_2,\mathfrak{m}_2)$ holds for every $u,v\in\mathbb S^{d-1}$ and thus it holds for the supremum over $u,v\in\mathbb S^{d-1}$. This supremum attains the operator norm, which implies that
\begin{equation}\label{eqn:LipDzgamma}
   \Vert D_z\gamma(z_2,\mathfrak{m}_1)-D_{z}\gamma(z_2,\mathfrak{m}_2)\Vert_\mathrm{op}\le 3R_1(1+2R_1R_z) \exp(2R_zR_1p^{-1}) W_p(\mathfrak{m}_1,\mathfrak{m}_2).
\end{equation}
 Automatically from Lemma \ref{lemma:2ndderivatives}, $D_z\gamma$ is $8R_1^3$-Lipschitz continuous in $z$ on $\bar B(R_1)\times \mathcal P(\bar B(R_1))$. Therefore, $D_z\gamma$ is Lipschitz continuous on $\bar B(R_1)\times \mathcal P(\bar B(R_1))$ with Lipschitz constant that is the maximum of $8R_1^3$ and the constant in equation \eqref{eqn:LipDzgamma}.

Now, consider finding the local Lipschitz constant of $\partial_\mu\gamma$. By the Mean Value Theorem on $\R$, for every $u,v\in \mathbb S^{d-1}$, there exists $t\in(0,1)$ with $z_t=z_1+t(z_2-z_1)$ such that
\[
u^T\left( \partial_\mu\gamma(z_1,\mathfrak{m}_1)(y_1)-\partial_\mu\gamma(z_2,\mathfrak{m}_1)(y_1)\right)v\le \left\Vert \sum_{i=1}^d (z_1-z_2)_i \partial_{z_i}\partial_\mu \gamma(z_t,\mathfrak{m}_1)(y_1) \right\Vert_\mathrm{op}.
\]
Hence, applying Lemma \ref{lemma:2ndderivatives} and then taking the supremum over $u,v \in \mathbb S^{d-1}$, we obtain
\begin{equation}\label{eqn:LipDmugamma3}
  \Vert \partial_\mu\gamma(z_1,\mathfrak{m}_1)(y_1)-\partial_\mu\gamma(z_2,\mathfrak{m}_1)(y_1)\Vert_\mathrm{op}\le 2R_1(2+3R_1R_z) \exp(2R_zR_1) \Vert z_1-z_2\Vert,   
\end{equation}
since for any $z_1,z_2\in \bar B(R_z)$ so too is $z_t\in \bar B(R_z)$. For the Lipschitz continuity of $\partial_\mu \gamma$ in measure and $y$, we use the expression for $\partial_\mu \gamma$ given in  \eqref{eqn:Dmugamma} and add and subtract suitable terms to get 
\begin{align}
    \Vert \partial_\mu\gamma(z_2,\mathfrak{m}_1)(y_1)-\partial_\mu\gamma(z_2,\mathfrak{m}_2)(y_2)\Vert_\mathrm{op}&\le \frac{\Vert z_2\Vert}{\gamma_2(z_2,\mathfrak{m}_1)}\left\Vert \exp(z_2^Ty_1)\,y_1-\exp(z_2^Ty_2)\,y_2\right\Vert\notag\\&+ \frac{\vert \exp(z_2^Ty_1)-\exp(z_2^Ty_2)\vert}{\gamma_2(z_2,\mathfrak{m}_1)}\notag\\&+\left\vert \frac{1}{\gamma_2(z_2,\mathfrak{m}_1)}-\frac{1}{\gamma_2(z_2,\mathfrak{m}_2)}\right\vert\exp(z_2^Ty_2)\left(1+2R_1R_z \right)\notag\\
    &+\frac{\exp(z_2^Ty_2)}{\gamma_2(z_2,\mathfrak{m}_2)}\Vert \gamma(z_2,\mathfrak{m}_1)-\gamma(z_2,\mathfrak{m}_2)\Vert \Vert z_2\Vert,\label{eqn:LipDmugamma}
\end{align}
Since $y_1,y_2\in \bar B(R_1)$, $z_2\in \bar B(R_z)$ and $\mathfrak{m}_1,\mathfrak{m}_2 \in \mathcal P(\bar B(R_1))$, the first term in the expression above is bounded via
\[
\frac{\Vert z_2\Vert}{\gamma_2(z_2,\mathfrak{m}_1)}\left\Vert \exp(z_2^Ty_1)\,y_1-\exp(z_2^Ty_2)\,y_2\right\Vert\le R_z\exp(2R_1R_z)\left(1+R_1R_z \right)\Vert y_1-y_2\Vert.
\]
Similarly, the second term is bounded by
\[
\frac{\vert \exp(z_2^Ty_1)-\exp(z_2^Ty_2)\vert}{\gamma_2(z_2,\mathfrak{m}_1)}\le \exp(2 R_1R_z)R_z \Vert y_1-y_2\Vert.
\]
For the third term, take any coupling $\bm{m}\in \Pi(\mathfrak{m}_1,\mathfrak{m}_2)$. By applying H\"older's inequality for any $p,q \in [1,\infty]$ with $p^{-1}+q^{-1}=1$, we obtain
\begin{align*}
    \left\vert \frac{1}{\gamma_2(z_2,\mathfrak{m}_1)}-\frac{1}{\gamma_2(z_2,\mathfrak{m}_2)}\right\vert&\le \frac{\int \exp(z_2^Ty_1)\vert 1-\exp(z_2^T(y_2-y_1))\vert\,d \bm{m}(y_1,y_2)}{\gamma_2(z_2,\mathfrak{m}_1)\gamma_2(z_2,\mathfrak{m}_2)}\\&\le\frac{\left(\int \exp(q z_2^Ty_1)d\mathfrak{m}_1\right)^{1/q}}{\gamma_2(z_2,\mathfrak{m}_1)} \frac{\exp(2R_zR_1)R_z}{\gamma(z_2,\mathfrak{m}_2)}\left( \int \Vert y_1-y_2\Vert^p d\bm{m}(y_1,y_2)\right)^{1/p}.
\end{align*}
Therefore, by taking the coupling $\bm{m}\in \Pi(\mathfrak{m}_1,\mathfrak{m}_2)$ that attains $W_p(\mathfrak{m}_1,\mathfrak{m}_2)$ yields
\begin{equation}\label{eqn:Lipinversegamma2}
    \left\vert \frac{1}{\gamma_2(z_2,\mathfrak{m}_1)}-\frac{1}{\gamma_2(z_2,\mathfrak{m}_2)}\right\vert\le R_z\exp((3+2p^{-1}) R_zR_1)W_p(\mathfrak{m}_1,\mathfrak{m}_2),
\end{equation}
Finally, by applying $\Lambda_{\mu,p}$-Lipschitz continuous of $\gamma$ with respect to its measure argument, Corollary \ref{corr:LipMu}, to the final term of \eqref{eqn:LipDmugamma} and the preceding estimates to bound the other terms of \eqref{eqn:LipDmugamma}, we deduce that there exists a constant $\Lambda_{\gamma\mu,p}'=\Lambda_{\gamma\mu,p}'(R_1,R_z)$ such that
\begin{equation}\label{eqn:LipDmugamma2}
    \Vert \partial_\mu\gamma(z_2,\mathfrak{m}_1)(y_1)-\partial_\mu\gamma(z_2,\mathfrak{m}_2)(y_2)\Vert_\mathrm{op}\le\Lambda_{\gamma\mu,p}' \left(\Vert y_1-y_2 \Vert + W_p(\mathfrak{m}_1,\mathfrak{m}_2)\right).
\end{equation}
Hence, combining bounds \eqref{eqn:LipDmugamma3} and \eqref{eqn:LipDmugamma2} yields that $\partial_\mu \gamma$ is jointly Lipschitz continuous on \\$\bar B(R_z)\times \mathcal P(\bar B(R_1)) \times \bar B(R_1)$.
\end{proof}

\end{document}
